% Generated draft: requires mathematical and source review.
\documentclass[11pt]{article}
\usepackage[margin=25mm]{geometry}
\usepackage{fontspec}
\setmainfont{Latin Modern Roman}
\usepackage{amsmath,amssymb,mathtools,unicode-math}
\setmathfont{Latin Modern Math}
\usepackage{enumitem,needspace,xurl,hyperref}
\setlength{\parindent}{0pt}
\setlength{\parskip}{5pt}
\setlength{\emergencystretch}{3em}
\newcommand{\N}{\mathbb{N}}
\newcommand{\Z}{\mathbb{Z}}
\newcommand{\Q}{\mathbb{Q}}
\newcommand{\R}{\mathbb{R}}
\newcommand{\C}{\mathbb{C}}
\newcommand{\F}{\mathbb{F}}
\newcommand{\A}{\mathbb{A}}
\newcommand{\PP}{\mathbb{P}}
\newcommand{\E}{\mathbb{E}}
\newcommand{\eps}{\varepsilon}
\newcommand{\dd}{\,\mathrm{d}}
\begin{document}
\section*{UnsolvedMath ID: 2307015}
\textbf{Source problem code:} catalogue entry
\subsection{Definitions and mathematical statement}
Let $D$ be a domain in the Riemann sphere with boundary of positive logarithmic capacity, and let $\omega_D^z$ denote harmonic measure viewed from $z\in D$, under an appropriate boundary measurability convention. A finite boundary point $p$ is angular if some $\eps>0$ has the following property: every component $V$ of $D\cap\{|w-p|<\eps\}$ with $p\in\partial V$ is contained in an angle of opening strictly less than $\pi$ and vertex $p$. The angle may depend on the component. The source defines angularity at infinity similarly without giving its coordinate convention. Write $A(D)$ for these points.

The source reports $\omega_D^z(A(D))=0$ for every $z\in D$ and asks for a direct proof, avoiding its cited unpublished Brownian-path theorem. It also says a nonempty $A(D)$ forces positive boundary capacity and that $A(D)$ can have positive linear measure. Its example takes a Cantor set $C\subset\{|z|=1\}$ and deletes the radial segments $\{rz:z\in C,\ 1/2\le r<1\}$ from the open unit disk, with the claimed conclusion $A(D)=C$. These are attributed contextual statements, not proved here.

For $1/2<\alpha<1$, the proposed replacement region is printed as
\[
Q_\alpha=\{x+iy:0<x<|y|^\alpha\}.
\]
The intended $B_\alpha(D)$ replaces the angles in the angular-point definition by translates of $Q_\alpha$. A literal finite-point version would require some $\eps>0$ such that every corresponding component $V$ is contained in $p+Q_\alpha$. The source does not specify rotations, dependence on components, normalization of the profile, or the version at infinity. The principal question is to determine which $\alpha\in(1/2,1)$ have the universal property
\[
\omega_D^z(B_\alpha(D))=0
\quad\text{for every domain }D\text{ in the stated capacity class and every }z\in D.
\]
This schema is conditional on recovering the exact intended replacement region and boundary-set convention.

The source says the analogous assertion is false at $\alpha=1/2$, using a disk. The displayed $Q_\alpha$ has two separated components and no point on $y=0$, so its literal geometry and this disk comparison require source review; no inequality reversal or alternative cusp is silently substituted. The quoted Brownian formula has a separate transcription defect: its displayed quantity is
\[
\left|\frac{\omega(t)-\omega(t_0)}{\omega(t)-\omega(t_0)}\right|,
\]
which, whenever defined, equals $1$ and cannot fill an open unit-circle arc of length $\pi$ as claimed. The intended directional expression is absent, and this formula cannot be used as a theorem.

A harder requested extension is to characterize the monotone profiles $f$ for which replacing the angles by translates of $\{x+iy:0<x<f(|y|)\}$ always yields a harmonic-measure-zero boundary set. Positivity, monotonicity direction, regularity, local normalization, and allowable coordinate changes are not stated. Analogous questions for Riemann surfaces and $n$-dimensional space are mentioned without further specifications. The supplied triage reports the item open as of 2018 with no definitive later resolution; no independent verification is supplied.
\subsection{Short English statement}
Determine which proposed cusp profiles force the corresponding boundary points to have harmonic measure zero, and seek a direct proof of the reported angular-point result. The printed Brownian expression is malformed, and the printed cusp's geometry needs reconciliation with its claimed disk endpoint example. Those defects prevent a unique corrected profile question from being inferred.
\subsection{Sources}
\begin{itemize}
\item Local source registry: \url{data/unsolvedmath/problems.json}, permanent UnsolvedMath ID $2307015$, source code catalogue entry, including the supplied OPEN-TRIAGE assessment. Dataset attribution: Ulam AI and UnsolvedMath Contributors; original-source attribution and rights are retained.
\item Hayman, \emph{Research Problems in Function Theory}, source-list date 2018, Problem 7.15, attributed to A. Dvoretzky: \url{https://arxiv.org/abs/1809.07200}. The metadata reports source-TeX extraction and July 29, 2026 access, with status and rights review flags.
\item The statement refers to an unpublished Brownian-path theorem, gives a Cantor radial-slit example and an endpoint disk claim, and supplies no separate authoritative citation or proof for them. The displayed Brownian formula is internally defective. No linked source was inspected for this draft.
\end{itemize}
\end{document}
